\chapter{Theory Bridge, Assumptions, and Non-Claims}
\label{ch:theory-bridge}

The monograph keeps theory and evidence coupled but not conflated.  The theory of
ORIUS proves that under explicit assumptions, degraded observation induces a hidden
safety problem and that a repair-and-certify layer can bound or eliminate the relevant
true-state violation event.  The empirical program then tests how much of that
theoretical surface survives domain transfer.

Two discipline rules follow.  First, theorem-grade claims remain strongest where the
theorem-to-code-to-artifact chain is deepest; in the current book that is still the
battery reference row, even though the reader-facing narrative is now universal-first.
Second, universality in the book means the architecture, runtime contract, and domain
chapter template are shared; it does not mean the book is allowed to flatten the parity
gate or hide unclosed rows.  This is why the monograph keeps one explicit equal-domain
parity matrix in view instead of letting narrative ambition substitute for closure.
